% 导言区：《近场动力学理论与数值方法》

% ---------- 页面设置 ----------
\usepackage[a4paper, top=2.5cm, bottom=2.5cm, left=3cm, right=2.5cm]{geometry}

% ---------- 常用宏包 ----------
\usepackage{amsmath}
\usepackage{amssymb}
\usepackage{amsthm}
\usepackage{graphicx}
\usepackage[font=small,labelsep=quad]{caption}
\graphicspath{{figures/}}
\usepackage{booktabs}
\usepackage{xcolor}
\usepackage{enumitem}
\usepackage{titlesec}
\usepackage{fancyhdr}
\usepackage{tikz}
\usepackage{listings}
\usepackage[colorlinks=true, linkcolor=blue, citecolor=blue, urlcolor=blue]{hyperref}

% ---------- 章节标题样式 ----------
\ctexset{
  chapter = {
    format = \centering\Huge\bfseries,
    name   = {第,章},
    number = \arabic{chapter},
  },
  section = {
    format = \Large\bfseries,
  },
  subsection = {
    format = \large\bfseries,
  },
}
% ---------- 目录深度设置 ----------
\setcounter{tocdepth}{1}
% Patch \addcontentsline to respect tocdepth at write time
% (titlesec writes all levels to .toc by design; this patch filters at write time)
\makeatletter
\def\pd@toclevel@section{1}
\def\pd@toclevel@subsection{2}
\def\pd@toclevel@subsubsection{3}
\def\pd@toclevel@paragraph{4}
\def\pd@toclevel@subparagraph{5}
\let\pd@addcontentsline@orig\addcontentsline
\renewcommand{\addcontentsline}[3]{%
  \def\pd@tmpa{#1}%
  \def\pd@tmpb{toc}%
  \ifx\pd@tmpa\pd@tmpb
    \@ifundefined{pd@toclevel@#2}{%
      \pd@addcontentsline@orig{#1}{#2}{#3}%
    }{%
      \ifnum\@nameuse{pd@toclevel@#2}>\c@tocdepth\relax\else
        \pd@addcontentsline@orig{#1}{#2}{#3}%
      \fi
    }%
  \else
    \pd@addcontentsline@orig{#1}{#2}{#3}%
  \fi
}
\makeatother
\captionsetup[figure]{name=图,labelsep=quad}
\captionsetup[table]{name=表,labelsep=quad}

% ---------- 页眉页脚 ----------
\pagestyle{fancy}
\fancyhf{}
\fancyhead[C]{\leftmark}
\fancyfoot[C]{\thepage}
\setlength{\headheight}{13.6pt}

% ---------- 定理环境 ----------
\newtheorem{theorem}{定理}[chapter]
\newtheorem{definition}{定义}[chapter]
\newtheorem{lemma}{引理}[chapter]
\newtheorem{remark}{注记}[chapter]
\newtheorem{example}{例}[chapter]

% ---------- 代码环境 ----------
\lstset{
  basicstyle       = \ttfamily\small,
  keywordstyle     = \color{blue},
  commentstyle     = \color{gray},
  stringstyle      = \color{red},
  frame            = single,
  breaklines       = true,
  numbers          = left,
  numberstyle      = \tiny\color{gray},
  showstringspaces = false,
}
